\subsubsection{Unitaries, Pauli Gates, and Eigenstates of Z, X, Y, H}\label{ex:unitaries-pauli-gates-and-eigenstates-of-z-x-y-h}

These definitions must become automatic before starting the exercises:
\begin{itemize}
    \item $Z$ gate (phase flip):
    \begin{equation*}
        Z = \begin{bmatrix}
            1 & 0 \\
            0 & -1
        \end{bmatrix}
    \end{equation*}
    \begin{equation*}
        Z \ket{0} = \begin{bmatrix}
            1 & 0 \\
            0 & -1
        \end{bmatrix} \begin{bmatrix}
            1 \\ 0
        \end{bmatrix} =
        \begin{bmatrix}
            1 \cdot 1 + 0 \cdot 0 \\
            0 \cdot 1 + (-1) \cdot 0
        \end{bmatrix} =
        \begin{bmatrix}
            1 \\ 0
        \end{bmatrix} =
        \ket{0}
    \end{equation*}
    \begin{equation*}
        Z \ket{1} = \begin{bmatrix}
            1 & 0 \\
            0 & -1
        \end{bmatrix} \begin{bmatrix}
            0 \\ 1
        \end{bmatrix} =
        \begin{bmatrix}
            1 \cdot 0 + 0 \cdot 1 \\
            0 \cdot 0 + (-1) \cdot 1
        \end{bmatrix} =
        \begin{bmatrix}
            0 \\ -1
        \end{bmatrix} =
        -\ket{1}
    \end{equation*}
    \item $X$ gate (bit flip):
    \begin{equation*}
        X = \begin{bmatrix}
            0 & 1 \\
            1 & 0
        \end{bmatrix}
    \end{equation*}
    \begin{equation*}
        X \ket{0} = \begin{bmatrix}
            0 & 1 \\
            1 & 0
        \end{bmatrix} \begin{bmatrix}
            1 \\ 0
        \end{bmatrix} =
        \begin{bmatrix}
            0 \cdot 1 + 1 \cdot 0 \\
            1 \cdot 1 + 0 \cdot 0
        \end{bmatrix} =
        \begin{bmatrix}
            0 \\ 1
        \end{bmatrix} =
        \ket{1}
    \end{equation*}
    \begin{equation*}
        X \ket{1} = \begin{bmatrix}
            0 & 1 \\
            1 & 0
        \end{bmatrix} \begin{bmatrix}
            0 \\ 1
        \end{bmatrix} =
        \begin{bmatrix}
            0 \cdot 0 + 1 \cdot 1 \\
            1 \cdot 0 + 0 \cdot 1
        \end{bmatrix} =
        \begin{bmatrix}
            1 \\ 0
        \end{bmatrix} =
        \ket{0}
    \end{equation*}
    \item $Y$ gate (bit and phase flip):
    \begin{equation*}
        Y = \begin{bmatrix}
            0 & -i \\
            i & 0
        \end{bmatrix}
    \end{equation*}
    \begin{equation*}
        Y \ket{0} = \begin{bmatrix}
            0 & -i \\
            i & 0
        \end{bmatrix} \begin{bmatrix}
            1 \\ 0
        \end{bmatrix} =
        \begin{bmatrix}
            0 \cdot 1 + (-i) \cdot 0 \\
            i \cdot 1 + 0 \cdot 0
        \end{bmatrix} =
        \begin{bmatrix}
            0 \\ i
        \end{bmatrix} =
        i\ket{1}
    \end{equation*}
    \begin{equation*}
        Y \ket{1} = \begin{bmatrix}
            0 & -i \\
            i & 0
        \end{bmatrix} \begin{bmatrix}
            0 \\ 1
        \end{bmatrix} =
        \begin{bmatrix}
            0 \cdot 0 + (-i) \cdot 1 \\
            i \cdot 0 + 0 \cdot 1
        \end{bmatrix} =
        \begin{bmatrix}
            -i \\ 0
        \end{bmatrix} =
        -i\ket{0}
    \end{equation*}
    \item $H$ gate (Hadamard):
    \begin{equation*}
        H = \frac{1}{\sqrt{2}}\begin{bmatrix}
            1 & 1 \\
            1 & -1
        \end{bmatrix}
    \end{equation*}
\end{itemize}
Also, a general rule for the action of a unitary gate $U$ on a qubit state $\ket{\psi}$:
\begin{equation*}
    U\left(\alpha \ket{0} + \beta\ket{1}\right)
\end{equation*}
Follow these four steps:
\begin{enumerate}
    \item Expand using linearity:
    \begin{equation*}
        U\left(\alpha \ket{0} + \beta\ket{1}\right) = \alpha U\ket{0} + \beta U\ket{1}
    \end{equation*}
    \item Apply the gate separately to each basis state.
    \item Replace $U\ket{0}$ and $U\ket{1}$ with their known values (from the definitions of the gates).
    \item Simplify.
\end{enumerate}
If a common factor such as $i$ appears everywhere, remember what we have learned in the previous section (\autopageref{sec:common-factors-global-phase-and-relative-phase}): it may simply be a \textbf{global phase} and can be ignored when calculating probabilities.

\highspace
Another central equation is:
\begin{equation*}
    U\ket{\psi} = \lambda \ket{\psi}
\end{equation*}
A state \ket{\psi} is an \textbf{eigenstate} of $U$ when applying $U$ does not change its direction: the result is only multiplied by a scalar $\lambda$, called the \textbf{eigenvalue}. For unitary gates, the eigenvalue has modulus one: $\left|\lambda\right| = 1$.

\highspace
\exercise Which is the eigenstate of the $Z$ gate? What is its eigenvalue?

\highspace
\solution Using the definition of the $Z$ gate, we have:
\begin{equation*}
    Z\ket{0} = \ket{0}
\end{equation*}
This means that $\ket{0}$ is an eigenstate of $Z$ with eigenvalue $\lambda = 1$. We also have:
\begin{equation*}
    Z\ket{1} = -\ket{1}
\end{equation*}
This means that $\ket{1}$ is also an eigenstate of $Z$, but with eigenvalue $\lambda = -1$. Thus, the $Z$ gate has two eigenstates: $\ket{0}$ and $\ket{1}$, with eigenvalues $1$ and $-1$, respectively.

\highspace
\exercise Which is the eigenstate of the $X$ gate? What is its eigenvalue?

\highspace
\solution Using the definition of the $X$ gate, we have:
\begin{equation*}
    X\ket{0} = \ket{1}
\end{equation*}
This means that $\ket{0}$ is \textbf{not} an eigenstate of $X$, because applying $X$ to $\ket{0}$ gives $\ket{1}$, which is a different state. Similarly:
\begin{equation*}
    X\ket{1} = \ket{0}
\end{equation*}
This means that $\ket{1}$ is also \textbf{not} an eigenstate of $X$. We can use $\ket{+}$ and $\ket{-}$, which are defined as:
\begin{equation*}
    \ket{+} = \frac{1}{\sqrt{2}}(\ket{0} + \ket{1}), \quad \ket{-} = \frac{1}{\sqrt{2}}(\ket{0} - \ket{1})
\end{equation*}
Let's check if these are eigenstates of the $X$ gate:
\begin{equation*}
    \begin{aligned}
        X\ket{+} &= X\left(\frac{1}{\sqrt{2}}(\ket{0} + \ket{1})\right) \\
        &= \frac{1}{\sqrt{2}}(X\ket{0} + X\ket{1}) \\
        &= \frac{1}{\sqrt{2}}(\ket{1} + \ket{0}) \\
        &= \frac{1}{\sqrt{2}}(\ket{0} + \ket{1}) \\
        &= \ket{+}
    \end{aligned}
\end{equation*}
Thus, $\ket{+}$ is an eigenstate of the $X$ gate with eigenvalue $\lambda = 1$. Let's check $\ket{-}$:
\begin{equation*}
    \begin{aligned}
        X\ket{-} &= X\left(\frac{1}{\sqrt{2}}(\ket{0} - \ket{1})\right) \\
        &= \frac{1}{\sqrt{2}}(X\ket{0} - X\ket{1}) \\
        &= \frac{1}{\sqrt{2}}(\ket{1} - \ket{0}) \\
        &= -\frac{1}{\sqrt{2}}(\ket{0} - \ket{1}) \\
        &= -\ket{-}
    \end{aligned}
\end{equation*}
Thus, $\ket{-}$ is also an eigenstate of the $X$ gate, but with eigenvalue $\lambda = -1$. Therefore, the $X$ gate has two eigenstates: $\ket{+}$ and $\ket{-}$, with eigenvalues $1$ and $-1$, respectively.

\highspace
\exercise Which is the eigenstate of the $Y$ gate? What is its eigenvalue?

\highspace
\solution Similarly, we can check the eigenstates of the $Y$ gate. Let's define the states:
\begin{equation*}
    \ket{+i} = \frac{1}{\sqrt{2}}(\ket{0} + i\ket{1}), \quad \ket{-i} = \frac{1}{\sqrt{2}}(\ket{0} - i\ket{1})
\end{equation*}
Let's check if these are eigenstates of the $Y$ gate:
\begin{equation*}
    \begin{aligned}
        Y\ket{+i} &= Y\left(\frac{1}{\sqrt{2}}(\ket{0} + i\ket{1})\right) \\
        &= \frac{1}{\sqrt{2}}(Y\ket{0} + iY\ket{1}) \\
        &= \frac{1}{\sqrt{2}}(i\ket{1} + i(-i\ket{0})) \\
        &= \frac{1}{\sqrt{2}}(i\ket{1} + \ket{0}) \\
        &= \frac{1}{\sqrt{2}}(\ket{0} + i\ket{1}) \\
        &= \ket{+i}
    \end{aligned}
\end{equation*}
Thus, $\ket{+i}$ is an eigenstate of the $Y$ gate with eigenvalue $\lambda = 1$. Let's check $\ket{-i}$:
\begin{equation*}
    \begin{aligned}
        Y\ket{-i} &= Y\left(\frac{1}{\sqrt{2}}(\ket{0} - i\ket{1})\right) \\
        &= \frac{1}{\sqrt{2}}(Y\ket{0} - iY\ket{1}) \\
        &= \frac{1}{\sqrt{2}}(i\ket{1} - i(-i\ket{0})) \\
        &= \frac{1}{\sqrt{2}}(i\ket{1} - \ket{0}) \\
        &= -\frac{1}{\sqrt{2}}(\ket{0} - i\ket{1}) \\
        &= -\ket{-i}
    \end{aligned}
\end{equation*}
Thus, $\ket{-i}$ is also an eigenstate of the $Y$ gate, but with eigenvalue $\lambda = -1$. Therefore, the $Y$ gate has two eigenstates: $\ket{+i}$ and $\ket{-i}$, with eigenvalues $1$ and $-1$, respectively.